\section{La escalabilidad del entrenamiento comienza siendo un problema de flujo de recursos}

"Diez mil GPUs" no significa copiar el programa de una sola tarjeta diez mil veces. Los parámetros del modelo, los gradientes, los estados del optimizador, las activaciones, los datos de entrada y los \textbf{checkpoint}s se mueven entre diferentes capas de almacenamiento y dispositivos; si cualquier flujo supera la capacidad del enlace, todo el entrenamiento se detendrá esperando. Por lo tanto, el hilo conductor de esta clase no es memorizar cinco nombres de \textit{parallelism}, sino preguntar repetidamente a cada estrategia: ¿qué se divide?, ¿qué se replica?, ¿cuál es el \textbf{comunicación colectiva} que restaura la equivalencia computacional y puede ocultarse con cómputo?

\subsection{¿Por qué los modelos grandes siguen creciendo en tamaño}

Las primeras tres diapositivas desglosan "escalabilidad" en tres ejes: modelo, datos y contexto. La primera página es un mapa reciente de modelos; la segunda alinea las tendencias de parámetros contra capacidad; y la tercera proporciona una orden de magnitud basada en la clase: billones de parámetros, unos $15$T tokens de entrenamiento y contextos de millones de tokens. No son una receta única, pero indican que la capacidad, el ancho de banda y el tiempo de iteración de una sola GPU están sujetos a presión simultánea.

\lecturefigure{slide-004.jpg}{Tendencias de lanzamiento de modelos: la escala del entrenamiento sigue expandiéndose}{CS25 V6 deck oficial, p.~4}

\lecturefigure{slide-005.jpg}{Contraste entre tendencias públicas de escala de modelos y cómputo de entrenamiento}{CS25 V6 deck oficial, p.~5}

\lecturefigure{slide-006.jpg}{Instantánea de orden de magnitud en clase: parámetros, tokens de entrenamiento y longitud de contexto}{CS25 V6 deck oficial, p.~6}

\begin{knowledgebox}{Lectura de la figura: los tres ejes de escala no son intercambiables}
El número de parámetros determina el estado base del modelo persistente y las operaciones matriciales por capa; el número de tokens de entrenamiento define el presupuesto total de iteraciones y la presión en la tubería de entrada; la longitud del contexto determina la activación por muestra, el conjunto de trabajo de atención y los micro-lotes disponibles. Aumentar DP solo permite consumir más lotes más rápido, no puede hacer que una secuencia excesivamente larga se ajuste automáticamente; aumentar CP puede dividir secuencias largas, pero no reduce el estado completo del modelo. ZeRO (\textit{Zero Redundancy Optimizer}, optimizador de redundancia cero) puede particionar el estado del optimizador, los gradientes y los parámetros por etapa, sin cambiar el número total de tokens de entrenamiento. Al analizar cada estrategia de \textit{parallelism} subsiguiente, primero devuélvala a estos tres ejes para determinar qué restricción realmente alivia y hacia dónde traslada los costos.
\end{knowledgebox}

\teachervoice{00:00:48--00:03:02, Tazi explica que esta clase se basa en el Ultra-Scale Playbook, pero incluye prácticas actualizadas de escalado MoE. Atribuye la presión simultáneamente a la lectura de datos, cada iteración, memoria y escritura del \textit{checkpoint}, en lugar de hablar solo del número de parámetros.}

\begin{warningbox}{La escala está relacionada con la inteligencia, pero no significa que la escala por sí sola cause inteligencia}
La clase usa modelos recientes para ilustrar que "entrenar más grande" sigue siendo una tendencia industrial, pero la correlación no aísla la calidad de los datos, la arquitectura, el entrenamiento posterior, el uso de herramientas y la contaminación de las evaluaciones. Lo verdaderamente necesario heredar del sistema son las órdenes de magnitud de recursos y las restricciones de ingeniería, no convertir alguna escala de parámetros en una ley de capacidad.
\end{warningbox}

\subsection{¿Qué se transporta realmente en una iteración de entrenamiento}

La infraestructura de entrenamiento debe leer continuamente el flujo de tokens, colocar los modelos y estados en la memoria del acelerador, mantener las activaciones entre forward/backward, ejecutar actualizaciones del optimizador y escribir periódicamente estados recuperables hacia el almacenamiento. Las magnitudes ilustradas en la figura a continuación—"~15T tokens read", "~1s iteration", "~80GB per GPU" y "~1TB checkpoint"—son solo ilustrativas; al interpretar la figura, considérelas como cuatro cuellos de botella independientes, no como una configuración precisa del mismo modelo.

\lecturefigure{slide-009.jpg}{Presión combinada del entrenamiento de LLM sobre datos, cómputo, memoria y E/S de \textit{checkpoint}}{CS25 V6 deck oficial, p.~9}

\begin{knowledgebox}{Lectura de la figura: considere un paso como una tubería multinivel}
Antes del paso de entrenamiento, el \textit{data loader} debe preparar muestras desde el almacenamiento y la memoria principal; luego la GPU lee parámetros y activaciones, ejecuta kernels y sincroniza mediante interconexión; al final del paso, el optimizador escribe los parámetros de vuelta, mientras que el ciclo del \textit{checkpoint} envía estados de recuperación más grandes al almacenamiento persistente. Cualquier segmento con insuficiencia de rendimiento generará un bloqueo (\textit{stall}) de forma diferente: la entrada insuficiente se manifiesta como tiempo libre antes del inicio de la GPU; la comunicación colectiva lenta se manifiesta como espera entre backward/forward; el \textit{checkpoint} lento se manifiesta como pausas largas periódicas. Al localizar el problema, primero determine a qué nivel pertenece el bloqueo, en lugar de atribuirlo uniformemente a "las GPUs no son lo suficientemente rápidas".
\end{knowledgebox}

\begin{knowledgebox}{Términos clave: primer grupo de objetos de recursos}
\textbf{HBM (High Bandwidth Memory, memoria de alto ancho de banda)} es DRAM de alto ancho de banda ubicada cerca de la GPU; a menudo se llama coloquialmente "memoria de video"; tiene capacidad limitada pero ofrece un ancho de banda muy alto para los kernels. \textbf{Optimizer state (estado del optimizador)} son estadísticas de actualización como el momento de primer orden $m$ y el de segundo orden $v$ que algoritmos como Adam/AdamW mantienen para cada parámetro; a menudo ocupan más bytes que los propios parámetros. \textbf{Activation (activación)} son resultados intermedios del forward, necesarios en backward para calcular gradientes. Los tres deben llevarse por separado.
\end{knowledgebox}

\subsection{Supuestos ocultos del baseline de una sola tarjeta y la acumulación de gradientes}

La lógica secuencial del entrenamiento en una sola tarjeta es clara: tome un lote, realice forward, calcule pérdida, backward para obtener gradientes, el optimizador lee parámetros y estado y escribe el modelo actualizado. Cuando el tamaño global del lote (GBS) es mucho mayor que el micro-lote que cabe en una sola tarjeta, se puede acumular gradientes mediante múltiples forward/backward antes de ejecutar una actualización.

\lecturefigure{slide-010.jpg}{Baseline forward--backward--optimizer de una sola GPU}{CS25 V6 deck oficial, p.~10}

\lecturefigure{slide-011.jpg}{Construcción de un lote global más grande mediante acumulación de gradientes}{CS25 V6 deck oficial, p.~11}

Sin embargo, la acumulación de gradientes solo resuelve "de cuántas muestras consta una actualización", no resuelve automáticamente si el modelo cabe o hace que los micro-lotes secuenciales sean más rápidos. La página siguiente escribe explícitamente estos dos pre-requisitos: el estado del entrenamiento debe poder residir primero en una sola tarjeta; si se espera a que todos los micro-lotes completen secuencialmente, tener más dispositivos no los aprovecha.

\lecturefigure{slide-012.jpg}{Dos pre-requisitos de acumulación en una sola tarjeta: el estado cabe y la espera secuencial es aceptable}{CS25 V6 deck oficial, p.~12}

Suponga que una actualización contiene $M$ micro-lotes, con pérdida $\ell_j(\theta)$ para cada micro-lote; entonces la acumulación de gradientes común se escribe como
\[
g=\frac{1}{M}\sum_{j=1}^{M}\nabla_\theta \ell_j(\theta),
\qquad
\theta' = \operatorname{Opt}(\theta,g,s).
\]
Donde $M$ es el número de pasos de acumulación, $\theta$ son los parámetros, $g$ es el gradiente promedio usado en esta actualización, $s$ es el estado del optimizador, $\operatorname{Opt}$ es la regla de actualización y $\theta'$ son los nuevos parámetros. Esta fórmula no prescribe que los $M$ gradientes deban generarse en el mismo dispositivo; el entrenamiento distribuido reorganiza exactamente la propiedad de todos estos cómputos y estados.

\subsection{Cinco ejes de división y marco unificado de evaluación}

El curso clasifica las estrategias en cinco categorías: Data, Tensor/Sequence, Pipeline, Context y Expert. Estas dividen a lo largo de los ejes batch, dimensión oculta, capa, secuencia y expert, respectivamente. El \textit{agenda} e introducción iniciales también enfatizan que estos mecanismos pueden usarse en GPUs o TPUs y no están limitados al pre-entrenamiento; desde dos dispositivos, la división de recursos puede tener sentido.

\lecturefigure{slide-013.jpg}{Ruta del curso: desde DP/ZeRO hasta TP, PP, CP y EP}{CS25 V6 deck oficial, p.~13}

Aquí aparece por primera vez el \term{ZeRO (\textit{Zero Redundancy Optimizer}, optimizador de redundancia cero)}, que no es un nuevo optimizador, sino un método de gestión de estados que particiona gradualmente el estado del optimizador, los gradientes y los parámetros, que originalmente se replicaban en los ranks de \textit{data-parallel}, por etapas. ZeRO-1/2/3 aumentan la variedad de estados particionados; los momentos de primer orden $m$ y segundo orden $v$ de Adam/AdamW siguen actualizándose según las fórmulas originales, cambiando solo la propiedad y el momento de recopilación.

\lecturefigure{slide-014.jpg}{Límites de aplicabilidad: múltiples aceleradores, etapas de entrenamiento y escala de dispositivos}{CS25 V6 deck oficial, p.~14}

El llamado \term{sharding (particionamiento)} consiste en cortar un objeto que originalmente se replicaba completamente en cada dispositivo en varios fragmentos (\textit{shards}), permitiendo que diferentes ranks posean solo una parte. Se debe explicar simultáneamente "qué se corta" y "cómo se corta a lo largo de qué grupo de dispositivos": la partición de parámetros, gradientes, lotes o secuencias no es lo mismo.

\lecturefigure{slide-015.jpg}{Objetivo uno: particionar modelo, gradientes y estado del optimizador}{CS25 V6 deck oficial, p.~15}

\lecturefigure{slide-016.jpg}{Objetivo dos: particionar datos y procesar micro-lotes en paralelo}{CS25 V6 deck oficial, p.~16}

\begin{importantbox}{Cuatro preguntas unificadas}
Al analizar cualquier estrategia de \textit{parallelism}, responda secuencialmente: primero, ¿qué dimensión de tensor o estado de entrenamiento se divide?; segundo, ¿qué objetos siguen siendo replicados?; tercero, ¿qué comunicación colectiva hace que los resultados sean matemáticamente equivalentes a un solo dispositivo?; cuarto, ¿esta comunicación está en la ruta crítica o puede superponerse con el cómputo? Si no puede aclarar estas cuatro preguntas, el llamado "ahorro de memoria" usualmente solo traslada costos a la red o al programador.
\end{importantbox}

\begin{knowledgebox}{Conceptos previos: el costo de comunicación no es solo el número de bytes}
El tiempo de una comunicación colectiva se puede escribir aproximadamente como $T\approx \alpha n_{\mathrm{steps}}+\beta V$: $\alpha$ representa la latencia de inicio y sincronización por ronda; $n_{\mathrm{steps}}$ depende de algoritmos como ring/tree y el número de ranks; $\beta$ es el tiempo de transmisión por byte y $V$ es el tráfico real. Los mensajes pequeños suelen dominarse por $\alpha$, mientras que los grandes son más sensibles al ancho de banda y la congestión topológica; los mensajes entre nodos también pasan por diferentes NICs, conmutadores y enrutadores. Por lo tanto, el mismo total de bytes, dividido en muchos \textit{all-gather} pequeños, puede ser más lento que pocos mensajes grandes; por el contrario, si un mensaje grande se inicia demasiado tarde, no se puede superponer. El diseño paralelo necesita optimizar simultáneamente tamaño del mensaje, frecuencia, grupos participantes y momento de inicio.
\end{knowledgebox}

\subsection{Resumen de la sección}

La entrada para la escalabilidad del entrenamiento no es "número de GPUs", sino un conjunto de restricciones de recursos: capacidad de estado, capacidad de activación, lote global, longitud de secuencia, estructura del modelo, topología de enlace y ancho de banda de almacenamiento. Las cinco estrategias de \textit{parallelism} reordenan la propiedad de estos recursos; los capítulos subsiguientes comenzarán con el más intuitivo \textit{data parallelism}, aumentando gradualmente la complejidad de comunicación y programación.


